\chapter{Ressources significatives}
\label{chap:ressources_significatives}

Ce chapitre rassemble cinq ressources transversales qui ont soutenu la progression du portfolio. Chaque ressource est présentée selon une logique de preuve : situation de mobilisation, outil réellement utilisé, résultat obtenu et trace vérifiable dans le dossier.

\section{Ressource 1 : Priorisation par critères explicites}

\textbf{Situation.} Plusieurs décisions techniques concurrentes devaient être traitées en parallèle, avec un risque de dispersion de l'équipe.

\textbf{Outil mobilisé.} Matrice d'Eisenhower, grille \textit{build vs buy} et hiérarchisation par priorité (HIGH/MEDIUM/LOW).

\textbf{Résultat observé.} Le débat d'opinions a été transformé en arbitrage structuré ; les sujets bloquants pour la certification ont été traités en premier.

\textbf{Trace.} Annexe~\ref{annex:ndt-sw-001} et Annexe~\ref{annex:backlog}.

\section{Ressource 2 : Agilité en contexte réglementé}

\textbf{Situation.} L'équipe devait itérer rapidement tout en maintenant la traçabilité attendue dans un environnement médical.

\textbf{Outil mobilisé.} Software Development Plan (SDP), \textit{Definition of Done} adaptée, articulation sprint/revue/rétrospective avec exigences documentaires.

\textbf{Résultat observé.} La conformité n'est plus traitée en fin de cycle ; elle devient un critère de conception intégré au flux de développement.

\textbf{Trace.} Annexe~\ref{annex:sdp-markdown} et section B4 de \textit{Situations professionnalisantes}.

\section{Ressource 3 : Écriture réflexive structurée}

\textbf{Situation.} Le risque principal du portfolio était de rester descriptif, sans démonstration claire de progression.

\textbf{Outil mobilisé.} Démarche réflexive structurée (problème, analyse, ressources, action, régulation), inspirée des cadres mobilisés en formation.

\textbf{Résultat observé.} Les chapitres passent d'un récit d'activités à une logique de transformation : avant/après, mécanismes causaux, limites et apprentissages.

\textbf{Trace.} Chapitre~\ref{chap:situations_professionnalisantes} et annexe du journal réflexif.

\section{Ressource 4 : Médiation et facilitation de convergence}

\textbf{Situation.} Les désaccords inter-profils (technique, qualité, usage) ralentissaient les décisions collectives.

\textbf{Outil mobilisé.} Recadrage, timeboxing, facilitation d'atelier, format Start/Stop/Continue en rétrospective.

\textbf{Résultat observé.} Les divergences ne sont plus subies ; elles sont traitées selon un protocole qui produit des décisions traçables et des actions fermables.

\textbf{Trace.} Annexe~\ref{annex:retro-b4} et section A2 de \textit{Situations professionnalisantes}.

\section{Ressource 5 : Culture de preuve qualité (IEC~62304)}

\textbf{Situation.} Les arbitrages de projet devaient être défendables dans la durée, au-delà de la mémoire des échanges.

\textbf{Outil mobilisé.} Notes de décision, backlog priorisé, matrices de traçabilité et contrôles de cohérence entre exigences, tests et livrables.

\textbf{Résultat observé.} Les décisions sont auditées plus facilement, les redites diminuent et la continuité entre décision et exécution est renforcée.

\textbf{Trace.} Annexes~\ref{annex:ndt-sw-001}, \ref{annex:backlog}, \ref{annex:c2-priorisation}.

\section{Synthèse}

La valeur de ces ressources ne réside pas dans leur addition, mais dans leur combinaison. Prioriser sans facilitation crée de nouveaux conflits ; itérer sans preuve crée de la dette ; réfléchir sans traces empêche d'apprendre collectivement. Leur articulation forme le socle méthodologique qui soutient l'ensemble du portfolio.
